\documentclass[../../main.tex]{subfiles}% 注意这里的文件路径不能用 ./main.tex ,否则用latexmk编译子文件会报错
\graphicspath{{\subfix{./image/}}} % 指定图片目录,后续可以直接使用图片文件名
% 注意这里的文件路径不能用 ../../image/ ,否则用latexmk编译子文件会报错

% 例如:
% \begin{figure}[H]
% \centering
% \includegraphics[scale=0.3]{图.png}
% \caption{}
% \label{figure:图}
% \end{figure}
% 注意:上述\label{}一定要放在\caption{}之后,否则引用图片序号会只会显示??.

\begin{document}

\section{主理想整环上的有限生成模}

总假定$D$为主理想整环（p.i.d.），零秩的自由$D$模为零模，秩为$m$的自由模$M$的秩记为$\mathrm{rank}\,M=r(M)=m$。

\begin{theorem}
设$M$是一个自由$D$模，则$M$的任一子模$N$也是自由$D$模，且$r(N)\leqslant r(M)$。
\end{theorem}
\begin{proof}

对$r(M)$使用数学归纳法。
当$r(M)=0$时，定理显然成立。

假设$r(M)=m-1$时定理成立，下证$r(M)=m$时结论成立。
取$M$的一组基$e_1,e_2,\cdots,e_m$，令
\begin{align*}
M'=\langle e_2,e_3,\cdots,e_m\rangle=\left\{\sum_{i=2}^{m}a_ie_i\mid a_i\in D,\,2\leqslant i\leqslant m\right\},
\end{align*}
则$M'$为$M$的秩为$m-1$的自由子模。

设$N\leqslant M$，分两种情况讨论：
\begin{enumerate}[(1)]
\item 若$N\subseteq M'$，则由归纳假设，$N$为自由$D$模，且$r(N)\leqslant r(M')<r(M)$。
\item 若$N\nsubseteq M'$，则存在$a=\sum_{i=1}^{m}a_ie_i\in N$满足$a_1\neq0$。
定义$D$的理想
\begin{align*}
I=\left\{a_1\in D\mid \sum_{i=1}^{m}a_ie_i\in N\right\},
\end{align*}
由$D$为主理想整环，故存在$d_1\in D\setminus\{0\}$使得$I=\langle d_1\rangle$，取$f_1=\sum_{i=1}^{m}d_ie_i\in N$。

任取$g=\sum_{i=1}^{m}a_ie_i\in N$，由$a_1\in I=\langle d_1\rangle$，存在$a_1'\in D$使得$a_1=a_1'd_1$，于是$g-a_1'f_1\in N\cap M'=N_1$，故$N=Df_1+N_1$。

若$a\in Df_1\cap N_1$，则存在$b_i\in D\,(2\leqslant i\leqslant m)$使得$\sum_{i=1}^{m}a_ie_i=\sum_{j=2}^{m}b_je_j$，故$a_1=0$。
又$a=a_0f_1$，故$a_1=a_0d_1$，结合$a_1=0$与$D$为整环得$a_0=0$，即$a=0$，因此$Df_1\cap N_1=\{0\}$，于是
\begin{align*}
N=Df_1\oplus N_1.
\end{align*}

$Df_1$是秩为1的自由$D$模，由$N_1\subseteq M'$及归纳假设，$N_1$为自由$D$模，因此$N$为自由$D$模，且
\begin{align*}
r(N)=1+r(N_1)\leqslant1+(m-1)=m.
\end{align*}
\end{enumerate}

\end{proof}

\begin{corollary}
主理想整环$D$上有限生成模的子模也是有限生成的。
\end{corollary}
\begin{proof}

设$D$模$M$的生成元为$g_1,g_2,\cdots,g_m$，$N$为$M$的子模。
秩为$m$的自由$D$模$D^{(m)}$以$e_1,e_2,\cdots,e_m$为基，存在模同态$\varphi:D^{(m)}\to M$满足$\varphi(e_i)=g_i,\,1\leqslant i\leqslant m$。

令$K=\varphi^{-1}(N)$，由同态基本定理，$K$为$D^{(m)}$的子模。
由前述定理，$K$为自由$D$模且$r(K)\leqslant m$。
设$f_1,f_2,\cdots,f_{r(K)}$为$K$的基，则$\varphi(f_1),\varphi(f_2),\cdots,\varphi(f_{r(K)})$为$N$的生成元，故$N$有限生成。

\end{proof}

\begin{remark}
一般环上自由模的子模不一定是自由模。
\end{remark}

\begin{example}
$\mathbb{Z}_6$是秩为1的自由$\mathbb{Z}_6$模，$\overline{1}$为其基。令$N=\{\overline{0},\overline{2},\overline{4}\}$，则$N$为$\mathbb{Z}_6$的非零子模。
由$|N|=3<6=|\mathbb{Z}_6|$，可知$N$不是自由$\mathbb{Z}_6$模。
\end{example}

\begin{definition}
设$M$是$R$模，$x\in M$。
\begin{enumerate}[(1)]
\item 若存在非零元$a\in R$使得$ax=0$，则称$x$为$R$上的\textbf{扭元}；否则称$x$为$R$上的\textbf{自由元（无关元）}。
\item 若$M$中每个元素都是扭元，则称$M$为\textbf{扭模}。
\item 若$M$中每个元素都是自由元，则称$M$为\textbf{无扭模}。
\end{enumerate}
\end{definition}

\begin{examples}
\begin{enumerate}[(1)]
\item 若$M$为有限Abel群，则$M$作为$\mathbb{Z}$模是扭模。
\item 设$M$为域$F$上的线性空间，则$M$作为$F$模是无扭模。
\item 设$V$是域$F$上的有限维线性空间，$\Phi$是$V$上的线性变换，$F[\lambda]$为$F$上一元多项式环，定义$f(\lambda)v=f(\Phi)v\,(f(\lambda)\in F[\lambda],\,v\in V)$，则$V$为$F[\lambda]$模，且$V$是扭模。
\item 整环$R$上的自由模一定是无扭模。
\item 有理数加群$\mathbb{Q}$作为$\mathbb{Z}$模是无扭模，但不是$\mathbb{Z}$上的自由模。
\end{enumerate}
\end{examples}

\begin{theorem}
设$M$是主理想整环$D$上的有限生成无扭模，则$M$是自由$D$模。
\end{theorem}
\begin{proof}

取$M$的一组生成元$x_1,x_2,\cdots,x_m$，则存在极大线性无关子集$\{x_1,x_2,\cdots,x_r\}$满足：
\begin{enumerate}[(1)]
\item $\sum_{i=1}^{r}a_ix_i=0$当且仅当$a_i=0\,(1\leqslant i\leqslant r)$，即$x_1,x_2,\cdots,x_r$线性无关；
\item 对任意$j>r$，存在不全为零的$a_{1j},a_{2j},\cdots,a_{rj},a_j\in D$使得
\begin{align}\label{eq:5.3.2}
\sum_{i=1}^{r}a_{ij}x_i+a_jx_j=0,\quad r<j\leqslant m.
\end{align}
\end{enumerate}

令$N=\langle x_1,x_2,\cdots,x_r\rangle=Dx_1\oplus Dx_2\oplus\cdots\oplus Dx_r$，则$N$为自由$D$模。
由条件，$a_j\neq0\,(r<j\leqslant m)$。

记$a=a_{r+1}a_{r+2}\cdots a_m$，定义映射$\eta:M\to M$满足$\eta(x)=ax,\,\forall x\in M$，则$\eta\in\mathrm{End}_D M$。
因$M$无扭，当$a\neq0$时，$ax=0$当且仅当$x=0$，故$\ker\eta=\{0\}$。

另一方面，
\begin{align*}
\eta(x_i)=ax_i\in N,\quad 1\leqslant i\leqslant r,
\end{align*}
\begin{align*}
\eta(x_j)=ax_j=\left(\prod_{k\neq j}a_k\right)a_jx_j=-\left(\prod_{k\neq j}a_k\right)\sum_{i=1}^{r}a_{ij}x_i\in N,\quad r+1\leqslant j\leqslant m.
\end{align*}
因此$\eta(M)\subseteq N$。
$N$为自由模，故$\eta(M)$也为自由模。结合$M\cong\eta(M)$，得$M$为自由$D$模。

\end{proof}

\begin{definition}
设$M$是$R$模，$x\in M$，定义$x$的\textbf{零化子}
\begin{align*}
\mathrm{ann}\,x=\{a\in R\mid ax=0\}.
\end{align*}
\end{definition}

\begin{remark}
$\mathrm{ann}\,x$是$R$的理想；若$x\neq0$，则$x$为扭元当且仅当$\mathrm{ann}\,x\neq\{0\}$。
\end{remark}

\begin{theorem}
设$R$为整环，$M$为$R$模，则$M$中扭元的集合
\begin{align*}
\mathrm{Tor}M=\{x\in M\mid \exists a\in R\setminus\{0\},\text{使得 }ax=0\}=\{x\in M\mid \mathrm{ann}\,x\neq\{0\}\}
\end{align*}
为$M$的子模，且商模$M/\mathrm{Tor}M$为无扭模。
\end{theorem}
\begin{proof}

任取$x,y\in\mathrm{Tor}M$，存在$a,b\in R\setminus\{0\}$使得$ax=by=0$，于是
\begin{align*}
ab(x-y)=b(ax)-a(by)=0,
\end{align*}
对任意$c\in R$，有
\begin{align*}
a(cx)=c(ax)=0.
\end{align*}
故$\mathrm{Tor}M$为$M$的子模。

记$\overline{M}=M/\mathrm{Tor}M$，任取$\overline{x}\in\overline{M}$，若存在$a\in R\setminus\{0\}$使得$a\overline{x}=0$，则$ax\in\mathrm{Tor}M$，故存在$b\in R\setminus\{0\}$使得$bax=0$，即$x\in\mathrm{Tor}M$，于是$\overline{x}=0$。
因此$\overline{M}$是无扭模。

\end{proof}

\begin{theorem}
设$D$为主理想整环，$M$为有限生成$D$模，则存在$M$的自由子模$N$使得
\begin{align}\label{eq:5.3.3}
M=\mathrm{Tor}M\oplus N,
\end{align}
且$N$在同构意义下唯一。
\end{theorem}
\begin{proof}

记$\overline{M}=M/\mathrm{Tor}M$，由前述定理，$\overline{M}$为无扭模。
设$x_1,x_2,\cdots,x_m$为$M$的生成元，$\pi:M\to\overline{M}$为自然同态，则$\{\overline{x}_i=\pi(x_i)\mid 1\leqslant i\leqslant m\}$为$\overline{M}$的生成元，故$\overline{M}$为有限生成无扭$D$模。
由前述定理，$\overline{M}$为自由$D$模。

取$e_1,e_2,\cdots,e_r\in M$使得$\overline{e}_1=\pi(e_1),\overline{e}_2=\pi(e_2),\cdots,\overline{e}_r=\pi(e_r)$为$\overline{M}$的基，令
\begin{align*}
N=De_1+De_2+\cdots+De_r.
\end{align*}

若存在$a_i\in D\,(1\leqslant i\leqslant r)$使得$\sum_{i=1}^{r}a_ie_i=0$，则
\begin{align*}
\pi\left(\sum_{i=1}^{r}a_ie_i\right)=\sum_{i=1}^{r}a_i\pi(e_i)=0,
\end{align*}
由$\overline{e}_1,\cdots,\overline{e}_r$线性无关得$a_1=a_2=\cdots=a_r=0$，故$e_1,e_2,\cdots,e_r$线性无关，于是
\begin{align*}
N=De_1\oplus De_2\oplus\cdots\oplus De_r,\quad N\cap\mathrm{Tor}M=\{0\},
\end{align*}
即式\eqref{eq:5.3.3}成立。

若另有分解$M=N'\oplus\mathrm{Tor}M$，则$N'\cong M/\mathrm{Tor}M\cong N$，故$r(N')=r(N)$，即$N$在同构意义下唯一。

\end{proof}

\begin{corollary}
若$M$是有限生成的$D$模，则$\mathrm{Tor}M$也是有限生成的$D$模。
\end{corollary}
\begin{proof}

该结论是主理想整环上有限生成模的子模有限生成的直接推论。

\end{proof}



\end{document}